\subsection{Application-route cost against public PyOptiX}
\label{sec:app-extension}

We separately evaluated the public application routes against competent public
PyOptiX implementations.  The registered denominator contains nine
project-authored application mappings.  Before the executable freeze, three of nine
applications (four operation units) had complete, output-equivalent PyOptiX
owners: one Particle Tracking transition step, RT-2A1 graph triangle counting,
and LibRTS point- and range-containment counts.  RTNN, X-HD, RT-DBSCAN, RayDB,
Spatial RayJoin, and RT-BarnesHut remained unexecuted because a competent
PyOptiX owner satisfying the registered output contract was not frozen.  We do
not remove those six applications from the denominator or replace them with
the two synthetic tasks evaluated above.

Both arms consumed the same frozen derived-input identity and returned the same
complete public result for each executed unit.  The PyOptiX arm owned the
public context, module, program groups, pipeline, SBT, GAS, device buffers,
launch, device continuation or reduction, synchronization, and result
materialization; it did not call a private RTDL native path.  Prebuilt PyOptiX
PTX was hash-bound before worker zero, whereas ordinary V4 Particle and
Triangle callback compilation remained inside registered complete preparation.
Each endpoint used eight paired blocks, four in each arm order, and a fresh
process for every arm and block.  Ratios are medians of eight within-block
V4/PyOptiX ratios.  The displayed arm times are diagnostic medians over the
eight arm-specific block medians and therefore need not divide to the displayed
ratio.  A ratio above one means V4 is slower.

\begin{table*}[t]
  \caption{Frozen first-batch application-route comparison on one NVIDIA RTX
  A4500 (CC 8.6).  Times are milliseconds.  ``Range'' is the minimum and
  maximum of the eight within-block V4/PyOptiX ratios.  All output checks
  passed; all twelve performance observations are adverse to V4.}
  \label{tab:app-pyoptix-results}
  \scriptsize
  \centering
  \begin{tabular}{@{}llrrrr@{}}
    \toprule
    Unit & Endpoint & V4 & PyOptiX & V4/PyOptiX & Range \\
    \midrule
    Particle Tracking & Setup + run & 24020.112 & 2405.894 & 10.208 & 9.559--10.676 \\
    Particle Tracking & First execute & 54.817 & 0.728 & 75.533 & 57.803--82.877 \\
    Particle Tracking & Prepared & 18.035 & 0.476 & 37.928 & 35.338--41.236 \\
    Triangle RT-2A1 & Setup + run & 11847.933 & 3508.170 & 3.411 & 3.292--3.591 \\
    Triangle RT-2A1 & First execute & 3516.714 & 1312.973 & 2.732 & 2.449--2.806 \\
    Triangle RT-2A1 & Prepared & 602.856 & 437.096 & 1.379 & 1.300--1.415 \\
    LibRTS point & Setup + run & 5502.395 & 4147.360 & 1.312 & 1.061--2.021 \\
    LibRTS point & First execute & 10.796 & 0.902 & 12.059 & 9.440--17.776 \\
    LibRTS point & Prepared & 8.916 & 0.230 & 38.467 & 34.687--41.240 \\
    LibRTS range & Setup + run & 4565.787 & 4173.109 & 1.080 & 1.009--1.196 \\
    LibRTS range & First execute & 12.352 & 0.890 & 13.422 & 12.182--23.731 \\
    LibRTS range & Prepared & 9.140 & 0.233 & 39.083 & 34.129--42.385 \\
    \bottomrule
  \end{tabular}
\end{table*}

The output obligations were a complete ordered $5000\mathbin{\times}3$ U32
Particle matrix, checked U64 graph count 2,224,385, checked U64 LibRTS point
count 112,729, and checked U64 LibRTS range count 105,826.  Triangle and LibRTS
therefore establish count routes, not relation materialization.  The registered
complete endpoint starts after a shared loader, then includes method import,
preparation, execution, required public checks and materialization, and close.
The excluded loader includes domain work: Particle uses already encoded arrays,
Triangle constructs its degree-oriented filtered and deduplicated CSR, and
LibRTS loads cached index columns while parsing WKT queries into points or MBRs.
This endpoint is method setup plus execution from a derived representation, not
raw-domain end-to-end application time; the table's Setup + run label therefore
means setup and run after shared preprocessing.  We report the protocol's
\texttt{first\_result} field as First execute after preparation: it is not cold
start or necessarily a pure launch because work inside the application execute
closure remains timed.  Execute timing also includes the registered synchronous
output/oracle and status checks, output digest, and compact evidence projection;
it is not isolated native-runtime latency.  Prepared measures one
complete public action on retained owners and retained registered input, after
one untimed warmup.

The route responsibilities are not uniform.  Particle authors still provide
mesh-to-face and adjacency encoding, strict-interior transition semantics,
queries, output meaning, and an oracle; RTDL compiles the supported four-role
callback and owns its typed ABI, OptiX wrapper, prepared lifecycle, result
checks, and traversal receipt.  This callback is supplied by the project
standard library rather than independently authored by a study participant.
The V4 wrapper enumerates any-hit candidates,
applies canonical $(t,\mathit{primitive\_id})$ selection, ignores each physical
intersection, and then invokes one logical closest-hit leaf.  PyOptiX instead
disables any-hit and uses native closest-hit.  Their frozen outputs match, but
their physical work is not identical.  Triangle authors select RT-2A1, construct
degree-oriented CSR segments and geometry, provide weights, combine segment
scalars, and check the graph result; RTDL executes the general Numba-leaf
device-column callback route, while CuPy performs checked-U64 weighted device
reduction.  Both Triangle arms rebuild per-segment geometry and GAS.  V4 also
builds and destroys its composed native module, program groups, pipeline, and
SBT per segment, while PyOptiX retains those objects from owner preparation.  The
PyOptiX device program is handwritten CUDA/OptiX compiled to the prebuilt PTX,
not Numba.  Thus the Triangle ratios include a lifecycle and program-generation
difference rather than isolating DSL code generation.  LibRTS authors select
point or range algebra and provide index and
query columns; RTDL verifies and executes a closed generic AABB-count
specialization rather than arbitrary callback lowering.

The closest steady result is Triangle prepared at 1.379$\times$.  The LibRTS
complete ratios, 1.312$\times$ and 1.080$\times$, are dominated by seconds of
preparation in both arms and do not imply prepared parity.  LibRTS prepared is
38.467--39.083$\times$ slower, and Particle prepared is 37.928$\times$ slower.
The final generic LibRTS device reduction downloads one U64 scalar rather than
a 100,000-element U32 count column, but it did not close the remaining physical
route gap.  Static inspection identifies broader V4 kernels, extra
synchronization, and route-specific transfer and materialization as candidate
mechanisms; this experiment does not causally allocate the total difference
among them.

These are registered post-loader route costs, not raw-domain application time
or intrinsic DSL overhead.  They do not support a speedup, near-PyOptiX
performance, broad application coverage, or reduced-authoring-effort claim.
The separate synthetic prepared measurements
in Table~\ref{tab:prepared} isolate different fixed tasks and cannot replace
this application evidence.  The complete 192-worker transaction retained zero
retries and discards and was independently recounted from its archived worker
records.
